\begingroup
\raggedright
\subsection{Fraction: Exact Rational Arithmetic}
C++17 (GCC/Clang). \texttt{Fraction(n,d)} represents $n/d$;
\texttt{Fraction(n)} is the integer $n$. Every result is reduced by
$\gcd$, its denominator is positive, and zero is stored as $0/1$.
Supports \texttt{+}, \texttt{-}, \texttt{*}, \texttt{/}, unary minus,
compound assignments and all six comparisons. Comparisons use exact
cross products, so the class also works with \texttt{sort},
\texttt{set} and \texttt{map} keys.

Numerator and denominator are signed 64-bit integers. Arithmetic
and cross products use \texttt{\_\_int128} intermediates, then reduce
before checking that the result fits; this is not arbitrary precision.
An unrepresentable reduced result throws \texttt{overflow\_error};
a zero constructor denominator throws \texttt{invalid\_argument},
and division by a zero fraction throws \texttt{domain\_error}.
Access parts with \texttt{num()} and \texttt{den()}.
Output uses \texttt{n/d}; an explicit \texttt{long double} conversion
is available when an approximation is wanted.
Each arithmetic operation costs $O(\log M)$ for gcd reduction;
comparisons use a constant number of integer operations.
Here $M$ bounds the intermediate integer magnitudes.
\lstinputlisting[style=C++,showstringspaces=false]{../mathematics/Fraction.hpp}
\begin{minipage}{\linewidth}
\begin{lstlisting}[style=C++,showstringspaces=false]
Fraction a(2, 4), b(1, 3); // a is reduced to 1/2.
std::cout << a << '\n';                  // 1/2
std::cout << a + b << ' ' << a - b << '\n'; // 5/6 1/6
std::cout << a * b << ' ' << a / b << '\n'; // 1/6 3/2
std::cout << (b < a) << '\n';             // 1
Fraction c(3, -6);                        // -1/2
c += 1;                                  // 1/2
\end{lstlisting}
\end{minipage}
\par
\endgroup
